\ifx\danslelivre\undefined
\documentclass[10pt]{../fichiers-configuration-latex/smfart}
\usepackage{ifthen}
\newboolean{english}
\newboolean{francais}
\setboolean{english}{false}
\setboolean{francais}{true}
\usepackage{palatino,euler}
\input{../fichiers-configuration-latex/paquets}
\selectlanguage{francais}
\input{../fichiers-configuration-latex/lemme-proposition-theoreme}
\input{../fichiers-configuration-latex/numerotation}
\input{../fichiers-configuration-latex/macros}
\input{../fichiers-configuration-latex/ucs_manquants}
\input{../.cv}

\externaldocument{../00-introduction/introduction}
\externaldocument{../01-quasi-excellents/excellent}
\externaldocument{../03-approximation/approximation}
\externaldocument{../04-Cohen-Gabber/Cohen-Gabber}
\externaldocument{../05-algebrisation-partielle/algebrisation}
\externaldocument{../06-log-lissite/log}
\externaldocument{../07-uniformisation-locale/uniformisation-faible}
\externaldocument{../08-modification-theorem/modification}
\externaldocument{../09-uniformisation-premiere-a-l/uniformisation-locale-premiere-a-l}
\externaldocument{../11-produits-orientes/produits-orientes}
\externaldocument{../12-descente-cohomologique-orientee/descente-orientee}
\externaldocument{../13-theoreme-de-finitude/finitude}
\externaldocument{../14-fonctions-dimension/fonctions-dimension}
\externaldocument{../15-Lefschetz-affine/Lefschetz-affine}
\externaldocument{../16-classes-de-Chern/gysin}
\externaldocument{../17-complexes-dualisants/dualite}
\externaldocument{../19-exemples-contre-exemples/contre-exemples}
\externaldocument{../20-rigidite/rigidite}
\externaldocument{../21-finitude-non-abelienne/finitude-non-abelienne}

\begin{document}

\begin{center}
II. Topologies adaptées à l'uniformisation locale\\
Fabrice Orgogozo
\end{center}

\version
\tableofcontents

\else
\setboolean{english}{false}
\setboolean{francais}{true}
\selectlanguage{francais}
\chapter{Topologies adaptées à l'uniformisation locale}
\label{2-topologies}
\addtocontents{toc}{par Fabrice Orgogozo}
\begin{center}
Fabrice Orgogozo
\end{center}
\fi

Dans cet exposé, $ℓ$ est un nombre premier, l'entier $1$
ou bien le symbole $∞$ et l'on note $ℓ′$ l'ensemble
des entiers naturels premiers à $ℓ$ où, par convention, $∞′=\{1\}$.

\section{Morphismes maximalement dominants et la catégorie $\alt ∕ S$}

\subsection{Morphismes maximalement dominants}

\subsubsection{}
Rappelons (\ega{IV}{1.1.4}) qu'un point d'un schéma est dit \emph{maximal} s'il est le
point générique d'une composante irréductible ou, de façon
équivalente, s'il est maximal pour l'ordre sur l'ensemble des
points du schéma défini par la relation : $x≥y$ si et seulement si $y$ est
une spécialisation de $x$ (c'est-à-dire si $y∈\sur{\{x\}}$).
Les points maximaux d'un schéma affine correspondent aux idéaux
premiers minimaux. Tout ouvert dense d'un schéma contient
la totalité des points maximaux.

\begin{définition2}\label{2-maximalement-dominant}
Un morphisme de schémas est dit \emph{maximalement dominant}\index{maximalement dominant}
s'il envoie tout point maximal de la source sur un point maximal du but.
\end{définition2}

Un morphisme entre schémas irréductibles est maximalement dominant si et seulement si il est
dominant. Il est clair que le composé de deux morphismes maximalement dominants est
maximalement dominant.

\begin{exemple2}
D'après \ega{IV₂}{2.3.4}, un morphisme plat, ou plus généralement quasi-plat (\emph{op. cit.},
2.3.3), est générisant (\ega{I′}{3.9.1}) donc maximalement dominant (\emph{op. cit.} 3.9.5).
\end{exemple2}

\begin{proposition2}\label{2-image maximalement dominant}
Soit $f:X→Y$ un morphisme maximalement dominant. Tout point maximal de $Y$ appartenant à l'image de $f$
est l'image d'un point maximal de $X$.
\end{proposition2}

\begin{démo}
Cela résulte du fait que $f$ est croissante pour le préordre ci-dessus
et du fait que tout point de $X$ a une générisation maximale.
\end{démo}

\begin{proposition2}\label{2-cb maximalement dominant par plat}
Soient $f:X→Y$ un morphisme maximalement dominant et $Y'→Y$ un morphisme
\emph{plat}. Alors, le morphisme $X'=X×_Y Y'→Y'$ est maximalement dominant.
\end{proposition2}

\begin{démo}
Cf. \ega{IV₂}{2.3.7} (ii).
\end{démo}

Rappelons la proposition suivante.

\begin{proposition2}[\ega{iv}{20.3.5}]
Soit $f:X→Y$ un morphisme maximalement dominant. Supposons que $X$
est \emph{réduit} et que $Y$ n'a qu'un nombre fini de composantes
irréductibles (par exemple, $Y$ nœthérien).
Alors, pour tout ouvert $U$ de $Y$ et tout ouvert $V$
dense dans $U$, l'ouvert $f^{-1}(V)$ est dense dans $f^{-1}(U)$.
\end{proposition2}

L'hypothèse sur $Y$ assure que tout point maximal de $U$ appartient à $V$.

\begin{proposition2}
\label{2-génériquement fini}
Soit $Y$ un schéma nœthérien et soit $f:X→Y$ un morphisme de type fini, \emph{maximalement
dominant}. Les conditions suivantes sont équivalentes :
\begin{enumerate}
\item Il existe un ouvert dense de $Y$ au-dessus duquel $f$ est fini.
\item Pour tout point maximal $η$ de $X$, l'extension $κ(η)/κ(f(η))$
est finie.
\end{enumerate}
\end{proposition2}

Un morphisme de schémas satisfaisant la condition (i)
ci-dessus est souvent dit \emph{génériquement fini} (en bas).

\begin{démo}
L'implication (i)⇒(ii) résulte du Nullstellensatz
et n'utilise pas les hypothèses de finitude faites ($Y$ nœthérien, $f$ de type fini).
Démontrons la réciproque. On peut supposer $Y$ irréductible et $X,Y$ réduits : on utilise le fait que $f$ est
fini si et seulement si $f_{\red}$ l'est. Par passage à la limite
(\ega{iv}{8.10.5.(x)}), on peut également supposer
que $Y$ est le spectre d'un corps $k$. On peut également
supposer $X$ irréductible donc intègre. Le résultat
est alors conséquence du fait qu'une algèbre intègre
entière sur un corps est un corps.
(Alternativement, on peut utiliser
l'inégalité $\dim(X)≤\dim(Y)=0$,
qui est un cas particulier de \ega{iv}{5.6.6}.)
\end{démo}

\begin{définition2}\label{2-L-maximalement fini}
On dit qu'un morphisme $f:X → Y$ est \emph{maximalement $ℓ′$-fini}
si pour tout point maximal $η$ de $X$, l'extension $κ(η)/κ(f(η))$
est finie et si pour tout point maximal $μ$ de $Y$ dans l'image de $f$,
il existe un point maximal $η$ de $X$ au-dessus de $μ$
tel que l'extension $κ(η)/κ(μ)$ soit de degré appartenant à $ℓ′$.
\end{définition2}

\subsubsection{}Lorsque $ℓ=1$, la seconde condition est vide.
Il est utile de faire les conventions de langage suivantes :
un morphisme maximalement $ℓ′$-fini et maximalement dominant
est dit \emph{maximalement $ℓ′$-fini dominant},
et un morphisme maximalement $1′$-fini « \emph{maximalement fini} ».

\begin{proposition2}\label{2-morphisme entre maximalement dominants gen finis}
Soient $S$ un schéma et $f:X→Y$ un $S$-morphisme entre $S$-schémas
maximalement dominants.
Si $Y/S$ est maximalement fini, $f$ est maximalement dominant.
Si l'on suppose de plus $X ∕ S$ maximalement fini, le morphisme $f$ est maximalement fini dominant.
\end{proposition2}


\begin{démo}
Soient $x$ un point maximal de $X$ et $s$ (resp. $y$) son image dans $S$
(resp. $Y$). Soit $y'≥y$ une générisation maximale de $y$. Les schémas
$X$ et $Y$ étant maximalement dominants sur $S$, $s$ est un point maximal et $y'$
est d'image $s$. Enfin, si $Y/S$ est maximalement fini, 
l'extension $κ(y')/κ(s)$ est finie de sorte que $y'$ est isolé dans la
fibre $Y_s$. Le point $y$, appartenant à l'adhérence de $y'$ dans
$Y_s$, coïncide donc avec $y'$ : le morphisme $f$ est maximalement dominant.
Si l'on suppose de plus $X/S$ maximalement fini,
la finitude de l'extension $κ(x)/κ(s)$ entraîne celle de l'extension
intermédiaire $κ(x)/κ(y)$.
\end{démo}

\subsection{La catégorie $\alt ∕ S$}

\subsubsection{}Soit $S$ un schéma nœthérien. Notons $η_S$ le schéma
coproduit (fini) de ses points maximaux.

\begin{définition2}
\label{2-alt}
On note $\alt/S$ \index{alt} la catégorie des $S$-schémas de type fini,
maximalement finis dominants, de source un schéma \emph{réduit}.
Les morphismes dans $\alt/S$ sont les $S$-morphismes.
\end{définition2}

\subsubsection{}Notons les faits suivants :
\begin{itemize}
\item le $S$-schéma $S_{\red}$ est final dans la catégorie $\alt /S$ ;
\item tout morphisme de $\alt/S$ est maximalement fini dominant
(\ref{2-morphisme entre maximalement dominants gen finis}) ;
\item les images inverses de diviseurs existent pour tout morphisme de $\alt/S$ (\ega{iv}{21.4.5.(iii)}) ;
\item si $X∈\Ob\,\alt/S$ et $S'→S$ est un morphisme \emph{réduit} (\ega{IV₂}{6.8.1}) avec $S'$
nœthérien, le produit fibré usuel $X'=X×_S S'$ est naturellement un objet de $\alt/S'$. Il
en est plus généralement ainsi de $X ′_\red$ dès lors que $S ′ → S$ est plat.
\end{itemize}

\begin{remarques2}
\begin{enumerate}
\item Le produit fibré usuel de deux $S$-schémas maximalement dominants
n'est pas nécessairement maximalement dominant, comme on peut
le constater lorsque $S=\Aff²$ et $X=Y$ sont l'éclatement en l'origine.
\item La définition originale de la catégorie $\alt ∕ S$, due à O. Gabber,
est moins restrictive sur le schéma $S$ :
il est supposé cohérent et ayant un nombre fini de composantes
irréductibles. Les objets de $\alt ∕ S$ sont alors les
$S$-schémas de type fini quasi-séparés réduits, maximalement dominants,
génériquement finis. Le cadre nœthérien semble suffisant pour
nos besoins. Signalons cependant que les « localisés »
d'un schéma nœthérien pour la topologie des altérations
(introduite ci-dessous) ne sont pas nécessairement nœthériens (cf. \ref{2-points de alt}).
\end{enumerate}
\end{remarques2}

\subsubsection{}
\label{2-dom}
Soit $X$ un $S$-schéma de type fini.
On note $X_\maxdom$ l'adhérence de l'image (ensembliste) de $X_{η_S}$
dans $X$, muni de la structure réduite. C'est la
réunion des composantes irréductibles de $X$ dominant une composante
irréductible de $S$, munie de la structure réduite.
Le foncteur $T↦T_\maxdom$ est adjoint à droite au foncteur
d'inclusion de la catégorie des schémas réduits, de type fini et
maximalement dominants sur $S$ dans la catégorie des $S$-schémas de type fini.

\begin{proposition2}\label{2-existence produit}
Les produits fibrés existent dans $\alt/S$.
\end{proposition2}

\begin{démo}
Soient $X→S'←Y$ deux flèches dans $\alt/S$ ; d'après
\ref{2-morphisme entre maximalement dominants gen finis}, les
schémas $X$ et $Y$ sont naturellement des objets de $\alt/S'$.
Le composé de deux morphismes maximalement finis dominants de type fini
étant de même nature, un produit de $X$ et $Y$, vus dans
$\alt/S'$, est — s'il existe — un produit fibré dans $\alt/S$.
On peut donc supposer $S=S ′$ et $S$ réduit.
Soient $X$ et $Y$ deux objets de $\alt/S$.
Il résulte formellement de l'existence du produit dans la catégorie des
$S$-schémas de type fini et de la propriété d'adjonction de $X↦ X_\maxdom$
que le schéma $(X×_S Y)_\maxdom$, muni des deux projections évidentes, est le produit
de $X$ et $Y$ dans la catégorie des schémas réduits, de type fini et
maximalement dominants sur $S$. Il appartient à $\Ob\,\alt/S$ car
$((X×_S Y)_\maxdom)_{η_S}=(X_{η_S}×_{η_S} Y_{η_S})_{\red}$ est fini sur $η_S$.
\end{démo}

\begin{proposition2}
\label{2-lprime-fini stable par cb}
La propriété d'être maximalement $ℓ′$-fini est stable par
changement de base.
\end{proposition2}

\begin{démo}
Il faut montrer que si $k ′ ∕ k$ est une extension finie de
degré premier à $ℓ$, où $k$ est de caractéristique $p$
et $K ∕ k$ est une extension quelconque
l'un des corps résiduels du produit tensoriel $K′=k'⊗_k K$ est de
degré premier à $ℓ$ sur $K$. On peut pour cela
supposer $k′ ∕ k$ étale ou radicielle.
Dans le premier cas, on écrit $K′= ∏_i K_i$, où $K_i$
est une extension étale de degré $d_i$ sur $K$, et on remarque
que si la somme $∑_i d_i=[k′:k]$ est première à $ℓ$,
il en est de même de l'un des $d_i$.
Le second cas ne se produit pas si $ℓ=p$ et est
trivial lorsque $ℓ≠p$, toute extension finie radicielle
étant de degré une puissance de $p$.
\end{démo}

\begin{proposition2}\label{2-adjoint preserve imm ouv, ps}
Soit $f:X→Y$ une immersion ouverte (resp. un morphisme propre et
surjectif, resp. quasi-fini) entre deux $S$-schémas de type fini. Le morphisme
$f_\maxdom:X_\maxdom→Y_\maxdom$ est une immersion ouverte (resp.
un morphisme propre et surjectif, resp. quasi-fini).
\end{proposition2}

\begin{démo}
Le cas d'une immersion ouverte est conséquence du fait général
suivant sur les topologies induites : la trace sur un \emph{ouvert} de l'adhérence d'une partie
coïncide avec l'adhérence de la trace de cette partie (cf. p. ex.
\bbk{TG}{I}{3}{1}{prop. 1} pour une variante).
Considérons maintenant le cas d'un morphisme propre et surjectif
$f:X→Y$. Par construction, les morphismes $X_\maxdom→X$ et
$Y_\maxdom→Y$ sont des immersions fermées ; les schémas
$X_\maxdom$ et $Y_\maxdom$ sont donc propres sur $Y$. Le $Y$-morphisme
$f_\maxdom$ est donc propre. Son image contient $f_{η_S}(X_{η_S})=Y_{η_S}$ donc son adhérence
$Y_\maxdom$. Le dernier cas est tout aussi trivial et laissé au lecteur.
\end{démo}

\subsubsection{}Notons que si $(U_i)_{i∈I}$ est un recouvrement
par des ouverts de Zariski d'un $S$-schéma de type fini $X$,
les ouverts $({U_i}_\maxdom)_{i∈I}$ recouvrent le schéma $X_\maxdom$.

\section{Topologies : définitions}

Dans ce paragraphe, on fixe un schéma nœthérien $S$.

\subsection{Topologie étale $ℓ′$-décomposée}

\subsubsection{}
\label{2-L-décomposé}
Nous dirons qu'un recouvrement étale $(X_i→ X)_{i∈I}$
d'un schéma $X$ est $ℓ′$-\emph{décomposé}
si tout point de $X$ peut être relevé en un point $x_i$ d'un $X_i$
tel que le degré $[κ(x_i): κ(x)]$ appartienne à $ℓ′$.
Lorsque $ℓ=1$ la condition imposée est vide et
l'on dit simplement que la famille constitue
un recouvrement étale. Lorsque $\mbox{ℓ=∞}$, on retrouve la définition de \cite[§1]{Nisnevic}
et l'on dit plutôt que le recouvrement est \emph{complètement} décomposé.
Il résulte de \ref{2-lprime-fini stable par cb} que
la propriété d'être un recouvrement étale $ℓ′$-décomposé est stable par
changement de base.

\subsubsection{}On appelle \emph{topologie étale $ℓ′$-décomposée}
la topologie de Grothendieck sur $\alt ∕ S$, notée $\et_{ℓ′}$,
définie par la prétopologie constituée des recouvrements
étales $ℓ′$-décomposés.

\subsection{Sorites sur le lieu $ℓ′$-décomposé}
\label{2-sorites lieu cd}

\subsubsection{}Pour chaque morphisme de schémas $f:Y → X$,
posons
\[
\mathrm{d\acute{e}c}_{ℓ′}(f) = \{x∈X:∃ y∈Y\textrm{ tel que } f(y)=x, [κ(x):κ(y)]\text{ fini appartenant à }ℓ′\}.
\]

Lorsque $ℓ=∞$, on retrouve l'ensemble $\mathrm{cd}(f)$ introduit
par Nisnevič. Nous utiliserons également cette notation.

\begin{proposition2}
Soit $f:Y → X$ un morphisme étale avec $X$ nœthérien.
L'ensemble $\mathrm{d\acute{e}c}_{ℓ′}(f)$ est
\emph{ind-constructible}, c'est-à-dire — $X$ étant nœthérien — réunion de parties localement fermées.
\end{proposition2}

\begin{démo}
On peut supposer $X$ et $Y$ intègres,
de points génériques notés $η$ et $μ$ respectivement.
Par récurrence nœthérienne, il suffit de montrer
que si $η$ appartient à l'ensemble $\mathrm{d\acute{e}c}_{ℓ′}(f)$,
celui-ci contient un ouvert de $X$.
On peut supposer de plus $X$ et $Y$ affines d'anneaux $A$ et $B$
respectivement, et le morphisme $A → B$ \emph{fini}.
La fonction $𝔭↦ \dim_{κ(𝔭)} B/𝔭B$, $X → 𝐍$,
est localement constante pour la topologie de Zariski
car $B$ est plat de présentation finie — donc localement
libre — sur $A$. Elle prend la valeur $[κ(μ): κ(η)]$ — première à $ℓ$ — en $η$.
Il en est donc ainsi au voisinage de $η$ ;
la conclusion en résulte aussitôt.
\end{démo}

Précisons un peu ce résultat lorsque $ℓ=∞$.

\begin{proposition2}\label{2-section cd}
Soit $f:Y→X$ un morphisme localement de type fini, avec $X$ nœthérien. Un point $x∈X$ appartient à
$\mathrm{cd}(f)$ si et seulement si il existe un sous-schéma $Z$ de $X$
contenant $x$ au-dessus duquel $f$ a une section.
\end{proposition2}

\begin{démo}
La condition est bien entendu suffisante.
Considérons $x∈\mathrm{cd}(f)$ ; c'est le point
générique du sous-schéma fermé réduit $\sur{\{x\}}$.
Par hypothèse, il existe une section au-dessus de ce point.
Le morphisme $f$ étant localement de présentation finie, cette section
s'étend par passage à la limite à un ouvert $Z=U⋂\sur{\{x\}}$
de $\sur{\{x\}}$, où $U$ est un ouvert de $X$.
\end{démo}

\begin{corollaire2}\label{2-cd noetherien}
Soient $X$ un schéma nœthérien et $(U_i\dessusdessous{f_i}{→}X)_{i∈I}$ un recouvrement de $X$ pour
la topologie étale $ℓ′$-décomposée. Il existe alors un sous-ensemble
fini $I₀⊂I$ tel que la famille $(U_i→X)_{i∈I₀}$ soit également couvrante.
\end{corollaire2}

Rappelons qu'un morphisme étale est, par définition, localement
de présentation finie.

\begin{démo}
D'après \ega{iv}{1.9.15}, l'espace topologique $X^{\mathrm{cons}}$,
dont l'espace sous-jacent est $X$ et dont les ouverts sont les parties
ind-constructibles de $X$, est \emph{compact}. Les ensembles $\mathrm{d\acute{e}c}_{ℓ′} (f_i)$
en constituent un recouvrement par des ouverts.
\end{démo}

\begin{proposition2}\label{2-cd limite}
Soit $(U_α\dessusdessous{f_α}{→}X_α)_{α∈A}$ un système projectif, filtrant, cartésien, de
morphismes étales entre schémas nœthériens, à morphismes de transition affines. Notons $f_∞:U_∞→X_∞$ le morphisme
induit sur la limite projective et, pour chaque $α$, $π_α$ la projection $X_∞ →  X_α$. 
On a 
$$
\mathrm{cd}(f_∞)=⋃_α π_α^{-1}\big(\mathrm{cd}(f_α)\big). 
$$
\end{proposition2}

\begin{démo}
L'inclusion du terme de droite dans le terme de gauche est évidente. 
Considérons réciproquement un point $x_∞$ dans $\mathrm{cd}(f_∞)$.
Le morphisme $f_∞$ a une section sur un sous-schéma de présentation
finie $Z_∞$ contenant $x_∞$. Le morphisme et la section se descendent par passage
à la limite à un niveau fini $α$ (cf. \ega{iv}{8.6.3, 8.8.2}).
\end{démo}

Les quatre énoncés précédents sont valables, \emph{mutatis mutandis},
lorsque l'on suppose simplement $X$ cohérent et le morphisme $f$
localement de présentation finie.
Pour référence ultérieure, signalons le lemme de descente
d'une section suivant.

\begin{proposition2}
\label{2-descente section}
Soient $k ′ ∕ k$ une extension finie de corps de degré premier à $ℓ$ et
$K ∕ k$ une extension finie de degré une puissance de $ℓ$. Posons $K
′= k′ ⊗_k K$. Si le morphisme $\Spec(K ′) → \Spec(K)$
possède une section, le morphisme $\Spec(k ′) → \Spec(k)$
possède également une section : c'est un isomorphisme.
\end{proposition2}

(Notons qu'un morphisme $X → \Spec(k)$ a une section
si et seulement si le morphisme $X_\red → \Spec(k)$
qui s'en déduit a une section.)

\begin{démo}
Cela résulte du fait que l'image de $k'$ dans $K$
par le morphisme composé $k ′ → K' → K$, où la seconde
flèche est la rétraction dont on suppose l'existence, est
à la fois de degré premier à $ℓ$ et de degré une puissance de $ℓ$ sur $k$.
\end{démo}

\subsection{Topologie des $ℓ′$-altérations}
\label{2-définition altL}

\subsubsection{}
On appelle \emph{topologie des $ℓ′$-altérations}
la topologie de Grothendieck sur $\alt ∕ S$, notée $\alt_{ℓ′}$,
définie par la prétopologie engendrée par
\begin{enumerate}
\item les recouvrements étales \emph{$ℓ′$-décomposés} ;
\item les morphismes propres \emph{maximalement $ℓ′$-finis} surjectifs.
\end{enumerate}

Prendre garde au fait que la seconde condition (« maximalement
$ℓ′$-fini ») porte sur les points maximaux tandis que la première
(« $ℓ′$-décomposé ») sur tous les points.

\begin{remarque2}
Les familles précédentes ne constituent
pas une prétopologie au sens de \sga{4}{II}{1.3} : la condition
de stabilité par composition n'est pas satisfaite. Les autres
conditions le sont, notamment la quarrabilité des morphismes (\ref{2-existence produit}).
\end{remarque2}

\subsubsection{}La topologie $\alt_{1′}$ est appelée \emph{topologie des
altérations}, notée simplement $\alt$.

\subsubsection{}
\label{2-stabilité par cb}
Notons pour référence ultérieure que si $(X_i → X)$ est une famille $\alt_{ℓ′}$-couvrante dans $\alt ∕ S$
et $S′ → S$ est un morphisme plat de source
un schéma localement nœthérien,
la famille $(X_i ×_S S′ → X×_S S′)$ de morphismes
dans $\alt ∕ S′$ est également $\alt_{ℓ′}$-couvrante
(cf. \ref{2-L-décomposé}, \ref{2-cb maximalement dominant par plat}).

\section{Formes standards}
Dans ce paragraphe, on fixe un schéma nœthérien $S$
et $X$ un objet de $\alt ∕ S$.
\subsection{Topologie étale}Le cas $ℓ=1$ de l'énoncé ci-dessous est un prototype
bien connu  des résultats que nous allons établir.

\begin{proposition2}
\label{2-forme standard étale}
Toute famille couvrante $(U_i → X)_{i ∈ I}$ pour la topologie étale
$ℓ′$-décomposée est dominée par une famille $\alt_{ℓ′}$-couvrante
du type
\[
(V_i → Y → X)_{i ∈ I}
\]
où $Y → X$ est fini, maximalement $ℓ′$-fini, surjectif et $(V_i → Y)_{i ∈ I}$ est un recouvrement
étale \emph{complètement} décomposé. Si $ℓ=1$, on peut supposer
que $(V_i → Y)_{i ∈ I}$ est un recouvrement par des ouverts de Zariski.
\end{proposition2}

\begin{démo}
On peut supposer l'ensemble $I$ fini (\ref{2-cd noetherien})
et $X$ intègre. Par passage à la limite, on peut
supposer de plus $X$ normal, de corps des fractions ayant
un groupe de Galois (absolu) pro-$ℓ$. (Le schéma $X$ n'est donc
pas nécessairement nœthérien.) Or, un morphisme étale $ℓ′$-décomposé
d'un tel schéma est nécessairement complètement décomposé
car le groupe de Galois des corps résiduels est également pro-$ℓ$.
Pour le complément lorsque $ℓ=1$, cf. p. ex. \cite[lemme 10.3]{Cycles@Orgogozo}.
\end{démo}

\subsubsection{}Il résulte immédiatement de la proposition
précédente que dans la définition du §\ref{2-définition altL},
on peut remplacer dans (i) la condition d'être
$ℓ′$-décomposé par celle d'être \emph{complètement}
décomposé (resp. de Zariski, si $ℓ=1$).

\subsection{Topologie des altérations}

Dans ce sous-paragraphe, on fixe un nombre premier $ℓ$.

\begin{théorème2}
\label{2-forme standard l'-altérations}
Supposons $X$ irréductible et quasi-excellent. Toute famille couvrante pour la topologie des $ℓ'$-altérations
de $X$ est dominée par une famille couvrante du type suivant :
\[
(V_i→ Y → X)_{i ∈ I},
\]
où $Y$ est un schéma intègre, $Y → X$ est propre et surjectif de degré
générique premier à $ℓ$ et $(V_i → Y)_{i ∈ I}$ est un recouvrement
pour la topologie complètement décomposée.
Si de plus $ℓ=1$, on peut supposer que $(V_i → Y)_{i ∈ I}$
est un recouvrement par des ouverts de Zariski.
\end{théorème2}

L'hypothèse de quasi-excellence est très certainement
superflue (procéder par passage à la limite).

Commençons par la démonstration du cas particulier $ℓ=1$,
bien qu'il résulte du cas général (joint
à \ref{2-forme standard étale} pour le complément).

\begin{démo}[Démonstration dans le cas où $ℓ=1$]
On peut supposer l'ensemble $I$ fini.
D'après \ref{2-forme standard étale}, il suffit de montrer
que si $(U_i→ X)_{i∈I}$ est un recouvrement de $X$ par des ouverts
de Zariski et $(X_i→ U_i)_{i∈I}$ une collection
de morphismes propres et surjectifs, il existe un morphisme
propre et surjectif $Y→X$ dans $\alt ∕ S$ et un
recouvrement par des ouverts
de Zariski $(V_i→Y)_{i ∈ I}$ tels que chaque morphisme
composé $V_i→ X$ se factorise à travers $X_i→ X$.
Choisissons pour chaque $i$ une compactification $\sur{X_i}→ X$
de $X_i → X$ ; on a $\sur{X_i}_{|U_i}=X_i$. Posons $Y=\sur{X_1}×_X \cdots
×_X \sur{X_r}$, où $I=\{1,…,r\}$, et $V_1=X_1×_X \sur{X_2}×_X\cdots
×_X \sur{X_r}$, $V_2=\sur{X_1}×_X X_2 ×_X \sur{X_3} ×_X \cdots ×_X
\sur{X_r}$, etc. Les ouverts $V_i$ recouvrent le schéma $\sur{X}$, qui
est propre et surjectif sur $X$. Par projection sur le $i$-ième facteur,
chaque $V_i$ s'envoie sur $X_i$. Quitte à appliquer le foncteur
$T↦ T_\maxdom$ (\ref{2-dom}), qui transforme un morphisme propre et surjectif (resp. un recouvrement
de Zariski) en un morphisme propre et surjectif (resp. en un recouvrement de
Zariski) (cf. \ref{2-adjoint preserve imm ouv, ps}), on obtient un recouvrement du type désiré
dans $\alt ∕ S$. Si $X$ est irréductible, on peut supposer $Y$ intègre.
\end{démo}

\begin{démo}[Démonstration dans le cas général]
Il suffit de vérifier que
si $(U_i→ X)_{i∈I}$ est un recouvrement étale complètement décomposé de $X$
et $(X_i→U_i)_{i∈I}$ une collection de morphismes propres,
surjectifs, maximalement $ℓ′$-finis, il existe une famille comme
dans l'énoncé la dominant. On peut supposer l'ensemble $I$ fini
et le schéma $X$ intègre et même normal compte-tenu
de l'hypothèse de quasi-excellence. Il est également loisible de supposer les $U_i$ connexes.
Notons qu'ils sont normaux. Observons d'autre part que l'on peut supposer les morphismes $X_i→U_i$ \emph{finis},
surjectifs, plats et de degré générique premier à $ℓ$ : chaque $X_i→X$ étant
génériquement plat ($X$ est réduit), quitte à remplacer $X$ par une
modification $X'→X$, on peut les platifier (\cite{Platification@Raynaud-Gruson},
 I 5.2.2), ce qui les rend \emph{finis}.
Quitte à ne considérer qu'une seule composante irréductible de chaque $X_i$,
de degré générique premier à $ℓ$ et à la normaliser,
on peut supposer les $X_i$ normaux connexes.
Résumons :

— les schémas $X$, $U_i$ et $X_i$ sont normaux intègres
de corps des fractions notés respectivement $K$, $K_i$
et $K_i′$ ;

— les morphismes $X_i → U_i$ sont finis de degrés
génériques premiers à $ℓ$.

Soit $L$ une clôture quasi-galoisienne sur $K$
d'une extension composée des $K_i′$,  et considérons
$X_L$ le normalisé de $X$ dans $L$. De même, considérons pour chaque $i$, le produit fibré
$U_{iL}=U_i×_X X_L$ (resp. le produit fibré
réduit normalisé $X_{iL}=(X_i×_X X_L)_{\red}^ν$). Compte tenu du
choix de $L$, pour chaque $i$ le morphisme $X_{iL}→U_{iL}$
a une section au-dessus de $\SP(L)$. Ce dernier étant fini
et $U_{iL}$ étant \emph{normal}, la section s'étend en une section $σ_i:U_{iL}→X_{iL}$.
Soit maintenant un $ℓ$-Sylow $S_ℓ$ de $\mathrm{Aut}(L/K)=G$
et notons $p$ la caractéristique du corps $K$. Si $ℓ=p$, les extensions
$K'_i/K$ sont donc étales de sorte que l'on peut supposer $L/K$ étale donc
\emph{galoisienne} ; l'extension $L^{S_ℓ}/K$ est alors de degré $(G:S_ℓ)$, premier à $ℓ$.
Si $ℓ≠p$, l'extension $L^{S_ℓ}/K$ est de degré $(G:S_ℓ)$ multiplié par une
puissance de $p$ ; c'est donc également un entier premier à $ℓ$.
Comme ci-dessus, notons $X_{L^{S_ℓ}}$ le normalisé de $X$ dans $L^{S_ℓ}$,
$U_{iL^{S_ℓ}}$ (resp. $X_{iL^{S_ℓ}}$) le produit fibré (resp. réduit normalisé)
de $U_i$ (resp. $X_i$) avec $X_{L^{S_ℓ}}$ au-dessus de $X$. Le morphisme
$X_{L^{S_ℓ}}→X$ est fini, de degré générique premier à $ℓ$.
D'après ce qui précède on a pour chaque $i$ un diagramme commutatif de 
schémas normaux :
$$
\xymatrix{
X_{iL^{S_\ell}} \ar[d]_{\text{premier à }\ell} & X_{iL} \ar[d] \ar[l] \\
U_{iL^{S_\ell}} & X_{iL} \ar[l]^{\text{puiss. de }\ell} \ar@/_1pc/[u]_{\sigma_i}
}
$$
En considérant isolément les composantes irréductibles du schéma
normal $U_{iL^{S_ℓ}}$ et, pour chacune d'entre elles, un point maximal de $X_{iL^{S_ℓ}}$ de degré
générique premier à $ℓ$ au-dessus, on montre immédiatement (cf. proposition \ref{2-descente section})
qu'il existe une section $U_{iL^{S_ℓ}}→X_{iL^{S_ℓ}}$.
Cela achève la démonstration de la proposition car
les $U_{iL^{S_ℓ}}$ forment un recouvrement pour la topologie complètement
décomposé du schéma $X_{L^{S_ℓ}}$, irréductible, de degré générique premier à $ℓ$ sur
$X$.
\end{démo}

\begin{remarques2}
\label{2-altération plus fine que étale cd}
\begin{enumerate}
\item Une famille couvrante pour la topologie des $ℓ ′$-altération le reste après changement de
base dans $\alt ∕ X$.
\item Il résulte du théorème que l'on obtient la même topologie si l'on remplace
la condition (i) de \ref{2-définition altL} par la condition
d'être un recouvrement étale \emph{complètement décomposé}.
Esquissons une preuve légèrement différente.
En passant à la limite sur les $Y$
comme dans le théorème précédent ($ℓ′$-altération de $X$)
on obtient un schéma normal dont le corps des fonctions a un
groupe de Galois absolu pro-$ℓ$. Cette propriété passe
aux corps résiduels (vérification triviale) si bien
qu'un revêtement étale $ℓ′$-décomposé est même
nécessairement \emph{complètement} décomposé.
\end{enumerate}
\end{remarques2}

Nous ferons également usage de la variante suivante
du théorème précédent.

\begin{théorème2}\label{2-forme standard pour annulation}
Supposons $X$ irréductible et quasi-excellent. Toute famille couvrante
pour la topologie des $ℓ'$-altérations de $X$
est dominée par une famille couvrante du type suivant :
\[
(W_i → V_i→ Y → X)_{i ∈ I},
\]
où tous les schémas sont irréductibles, $Y→X$ est propre birationnel,
$(V_i→Y)_{i ∈ I}$ est un recouvrement pour la topologie
complètement décomposée, et les morphismes
$W_i→V_i$ sont finis, plats, de degré premier à $ℓ$.
\end{théorème2}

Un énoncé semblable est également valable si $X$ n'est pas
irréductible : considérer le coproduit de ses composantes irréductibles.

\begin{démo}[Esquisse de démonstration]
Par platification, il suffit de montrer le résultat
d'échange suivant : si $Y → X$ est fini, plat, de degré
générique premier à $ℓ$ et $(V_i → Y)_{i ∈ I}$ est un
recouvrement étale complètement décomposé, il existe un
recouvrement étale complètement décomposé $(U_j → X)_{j ∈ J}$
et des morphismes $Z_j → W_j$, finis, plats, de degré
premier à $ℓ$ tels que la famille de morphismes
composés $(Z_j → X)$ domine celle des $(V_i → X)$.
Par passage à la limite, on peut supposer $X$ hensélien.
Par hypothèse, il existe une composante connexe $Y⁰$ de $Y$
qui est plate de degré premier à $ℓ$ sur $X$. Le schéma $Y⁰$ étant local hensélien,
il existe un indice $i$ tel que la restriction du
morphisme $V_i→Y$ à $Y⁰$ ait une section,
de sorte que le morphisme composé $Y⁰→X$ se factorise
à travers $V_i$. Cela permet de conclure.
\end{démo}

\section{Applications}


\subsection{Sorites}

\begin{proposition2}\label{2-voisinages fibre restent couvrants}
Soient $X$ un schéma nœthérien, $x$ un point de $X$,
$(X_i→X)_{i=1,\dots,n}$ un recouvrement pour la topologie des altérations
et, pour chaque indice $i$, un ouvert $X⁰_i↪X_i$ contenant la fibre
$(X_i)_x$. Il existe un voisinage ouvert de Zariski $U$ de $x$
tels que la famille $({X⁰_i}_{|U}→U)_i$ soit $\alt$-couvrante.
\end{proposition2}

Un cas particulièrement utile — et auquel on pourrait se
ramener d'après la proposition suivante  — est celui où
$X$ est local de point fermé $x$, de sorte que $U=X$.

\begin{démo}
D'après le théorème \ref{2-forme standard l'-altérations},
la famille $(X_i→X)_i$ est dominée
par une famille $(V_j↪Y→X)_j$ où $f:Y→X$ est notamment
propre et surjectif, et les $(V_j→Y)_j$ sont un recouvrement ouvert de $Y$. Soit $Y⁰_j$ l'image
inverse de $X⁰_i↪X_i$ par une factorisation $V_j→X_i$.
Par hypothèse, au-dessus du point $x$ de $X$, $Y⁰_j$ et $V_j$
coïncident, de sorte que $(Y⁰_j)_j$ est une famille d'ouverts
de $Y$ recouvrant la fibre $Y_x$. Leur réunion $Y⁰$ est un ouvert
de $Y$, contenant cette fibre, et on vérifie aussitôt que l'ouvert $U=X-f(Y-Y⁰)$
de $X$ convient.
\end{démo}

\begin{proposition2}\label{2-famille alt-couvrante et passage à la limite}
Soit $(X_α)_{α∈A}$ un système projectif filtrant de schémas
nœthériens affines à morphismes de transitions dominants. 
On suppose que la limite $X=\lim_α X_α$ est un schéma nœthérien
irréductible sur lequel un nombre premier $ℓ$ est inversible. 
Alors, toute famille $\altl$-couvrante de $X$ est dominée par l'image
inverse d'une famille $\altl$-couvrante de l'un des $X_α$.
\end{proposition2}

\begin{démo}
D'après le théorème \ref{2-forme standard l'-altérations}, joint au fait que les
diagrammes de morphismes se descendent, il suffit de démontrer le théorème
dans les cas particuliers suivants : la famille couvrante est
constituée d'un unique morphisme propre surjectif de degré générique
premier à $ℓ$ ou bien elle est finie, couvrante pour la
topologie étale complètement décomposée.
Traitons le premier cas. D'après \ega{iv}{8.10.5} (vi, xii),
tout morphisme propre et surjectif $Y→X$ se descend en un
morphisme propre et surjectif $Y_α→X_α$ ($α$ suffisamment grand).
Vérifions que l'on peut supposer $Y_α∈\Ob\,\alt/X_α$. D'après \ega{iv}{8.4.2}, on peut supposer
$X_α$ irréductible et $X→X_α$ maximalement dominant. Il en résulte,
puisque la flèche $Y_\maxdom → Y$ est un isomorphisme, que le morphisme $Y→Y_α$ se factorise à
travers $(Y_α)_\maxdom↪Y_α$. On peut donc supposer le
morphisme $Y_α → X_α$ maximalement dominant. Soit $η_α$ le
point générique de $X_α$. On doit vérifier que
la fibre générique $(Y_α)_{η_α}$ est intègre, finie de degré premier
à $ℓ$ sur $η_α$. Cela résulte du fait qu'elle l'est après
changement de base par $η→η_α$, où $η$ est le point
générique de $X$. Traitons maintenant le cas d'un recouvrement étale, complètement décomposé
$(U_i→X)_{i∈I}$ ($I$ fini). D'après \ega{iv}{17.7.8} (ii)
et \ega{iv}{8.10.5} (vi), cette famille provient d'une famille
$(U_{iα}→X_α)_{i∈I}$ couvrante pour la topologie étale d'un $X_α$,
pour $α$ suffisamment grand. D'après \ref{2-cd limite}, si pour 
tout $β≥α$ on note $f_β$ le morphisme $X_β×_{X_α} ∐_i U_{iα}→X_β$ et $π_β$
le morphisme $X→X_β$, on a $X=⋃_{β≥α} π_β^{-1}\big(\cd(f_β)\big)$.
C'est une réunion croissante d'ouverts du compact $X^\cons$ ;
pour $β$ assez grand on a donc $X=π_β^{-1}\big(\cd(f_β)\big)$.
Il résulte de \ega{iv}{8.3.11} que $\cd(f_β)=X_β$ pour $β$ assez
grand : sur $X_β$ le recouvrement étale est donc complètement
décomposé.
\end{démo}

\subsection{Caractérisation ponctuelle}

Terminons par une caractérisation de la topologie
des altérations, semblable à la caractérisation de 
la topologie étale à l'aide des anneaux
locaux strictement henséliens.
% cf. SGA 4 IV, 6.5.1 : points caractérisent (souvent) la topologie

\begin{théorème2}[\cite{Cohomological@Goodwillie-Lichtenbaum}, 3.5]\label{2-points de alt}
Soient $X$ un schéma nœthérien et $Y$ un objet de $\alt/X$.
La flèche $Y→X$ est couvrante pour la topologie des altérations 
si et seulement si elle est valuativement surjective au sens suivant : 
tout morphisme maximalement dominant $\SP(V)→X$, où $V$ est un anneau de valuation à corps
résiduel algébriquement clos se relève en un morphisme
$\Spec(V)→Y$.
\end{théorème2}

\begin{remarque2}
On dispose d'un analogue du théorème \ref{2-points de alt}, où l'on remplace le point générique géométrique
$\sur{η}$ (resp. la condition que le corps des fractions est algébriquement clos) par
le quotient $\sur{η}/H$ où $H$ est un $ℓ$-Sylow de $\Aut(κ(\sur{η})/κ(η))$ (resp. la condition
que le corps des fractions est parfait de groupe de Galois absolu un
pro-$ℓ$-groupe).
\end{remarque2}

\begin{démo}
Montrons que la condition est nécessaire.
D'après le théorème précédent, on peut supposer $Y→X$ propre et
surjectif. (Le cas où $Y$ est associé à un recouvrement par des
ouverts de Zariski est trivial car $\Spec(V)$ est local.)
Notons $η_V$ le point générique de $\Spec(V)$. Le morphisme
$η_V→X$ se relève en un morphisme (non unique) $η_V→Y$ 
d'après le Nullstellensatz, car $Y→X$ est surjectif 
et $κ(η_V)$ algébriquement clos.
Il résulte du critère valuatif de propreté que le morphisme
$η_V→Y$ s'étend en un $X$-morphisme $\Spec(V)→Y$.
(Remarquons qu'il n'est pas nécessaire de supposer $\Spec(V)→X$
maximalement dominant.)

Montrons que la condition est suffisante. 
On peut supposer pour simplifier $X$ affine intègre, de point
générique noté $η$. On peut également supposer $Y$ affine, 
de sorte qu'il existe une immersion ouverte $Y↪\sur{Y}$ dans un 
schéma propre et surjectif sur $Y$. 
Choisissons enfin un point générique géométrique $\sur{η}→X$ et 
considérons l'espace de Zariski-Riemann 
$\ZR_{\sur{η}}(X)$, limite des espaces annelés $X'$, où $X'$ est un schéma intègre,
propre et surjectif, muni d'un $X$-morphisme $\sur{η}→X'$.
On peut montrer qu'il est quasi-compact (cf. \emph{op. cit.} ou 
\cite{CA@Zariski-Samuel}, chap. VI, th. 40 pour une variante)
et que si $X=\Spec(A)$, l'application qui à un anneau de valuation $A⊆V⊆κ(\sur{η})$
associe le point de $\ZR_{\sur{η}}(X)$ correspondant par le critère
valuatif de propreté est une bijection. 

Tout relèvement $r:\sur{η}→Y$ du point générique géométrique de 
$X$ induit un $X$-morphisme $\sur{η}→\sur{Y}$ donc un morphisme
continu $π_r:\ZR_{\sur{η}}(X)→\sur{Y}$, qui se factorise à travers la composante
irréductible de $\sur{Y}$ atteinte par $r$. Par hypothèse,
les ouverts $π_r^{-1}(Y)$, pour $r$ variable, recouvrent 
$\ZR_{\sur{η}}(X)$. Par quasi-compacité, il existe donc un nombre 
fini de relèvements $r₁,\dots,r_n:\sur{η}→Y$
tels que $\ZR_{\sur{η}}→X$ se factorise à travers
la réunion des $n$ \emph{ouverts} images inverses de $Y$
dans le $X$-schéma propre et surjectif $\sur{Z}=\sur{Y}×_X 
\sur{Y}×\cdots×_X \sur{Y}$ ($n$ fois).
Par définition de l'espace de Zariski-Riemann, cette factorisation 
entraîne que $\sur{Z}$ est la \emph{réunion} de ces ouverts.
Ainsi, $Y→X$ peut-être raffiné en un recouvrement par des ouverts
de Zariski d'un schéma propre et surjectif. CQFD. 
\end{démo}

\subsection{Réduction des théorèmes d'uniformisation locale au cas hensélien}

Rappelons qu'un des objectifs de ce livre est de démontrer
les théorèmes d'uniformisation locale \refext{VII}{7-uniformisation}
et \refext{0}{0-uniformisation l'} dont nous reproduisons
l'énoncé.

\begin{théorème2}[\cite{Conference-Deligne@Gabber}, 1.1]\label{2-théorème uniformisation faible}
Soient $X$ un schéma nœthérien quasi-excellent et $Z$ un
fermé rare de $X$.
Il existe une famille finie de morphismes $(X_i→X)_{i∈I}$, couvrante pour la topologie des
altérations et telle que pour tout $i∈I$ on ait :
\begin{enumerate}
\item le schéma $X_i$ est régulier et connexe ;
\item l'image inverse de $Z$ dans $X_i$ est le support
d'un diviseur à croisements normaux stricts.
\end{enumerate}
\end{théorème2}

Par convention, l'ensemble vide est considéré comme
un diviseur strictement à croisements normaux : c'est
une somme indexée par l'ensemble vide.

\begin{théorème2}[\emph{op. cit.}, 1.3]\label{2-théorème uniformisation première à l}
Soient $X$ un schéma nœthérien quasi-excellent, $Z$ un
fermé rare de $X$ et $ℓ$ un nombre premier inversible sur $X$.
Il existe une famille finie de morphismes $(X_i\ra X)_{i∈I}$ , 
couvrante pour la topologie des $ℓ'$-altérations et telle que pour
tout $i\in I$ on ait :
\begin{enumerate}
\item le schéma $X_i$ est régulier et connexe ;
\item l'image inverse de $Z$ dans $X_i$ est le support d'un diviseur à croisements normaux stricts.
\end{enumerate}
\end{théorème2}

Dans la fin de cet exposé, nous allons nous contenter de
démontrer le fait suivant.

\begin{proposition2}\label{2-réduction cas hensélien}
Si l'un des théorèmes d'uniformisation est vrai pour tout schéma
$X$ local nœthérien hensélien excellent normal (resp. pour tout
schéma $X$ local nœthérien hensélien excellent normal de dimension au plus
un entier $d$ fixé), il est vrai en général (resp. pour tout schéma nœthérien
excellent de dimension finie inférieure ou égale à $d$).
\end{proposition2}

La réduction au cas où $X$ est local nœthérien \emph{complet}
est bien plus délicate ; elle fait l'objet de l'exposé suivant.

\begin{démo}
Supposons le théorème \ref{2-théorème uniformisation faible}
(resp. \ref{2-théorème uniformisation première à l})
démontré dans le cas local hensélien excellent.
Soit $X$ un schéma nœthérien quasi-excellent et $Z$ un fermé rare.
On peut supposer $X$ normal intègre car le morphisme de normalisation est couvrant
pour la topologie des $ℓ'$-altérations et l'image inverse de $Z$ reste
rare. Fixons $x∈X$. D'après \ega{iv}{18.7.6} l'hensélisé
$X_{(x)}$ de $X$ en $x$ est excellent et $Z_{(x)}$ est un fermé
rare de $X_{(x)}$. Il existe donc une famille finie 
de diagrammes
$$
\xymatrix{
Y \ar[d] & Y_i \ar[d] \ar[l] \\
X_{(x)} & X_i \ar[l]
}
$$
dans $\alt/X_{(x)}$, où $Y$ est intègre, propre et surjectif 
(resp. et de degré générique premier à $ℓ$) sur $X$, 
$(Y_i→Y)$ est une famille couvrante pour la topologie
de Zariski (resp. complètement décomposée) et la famille $(X_i→X)$
satisfait les conditions (i-ii).
Il résulte de la démonstration de la proposition \ref{2-famille alt-couvrante et passage à la limite}
que cette famille de diagrammes s'étend en
famille du même type sur un voisinage étale complètement
décomposé $U$ de $x$ dans $X$.
Il reste à vérifier que les propriétés (i) et (ii) sont bien
conservées. Si un morphisme $T→X_{(x)}$ de type fini,
avec $T$ \emph{régulier}, est le changement de base d'un morphisme
$V→U$ de type fini où $U$ est un voisinage étale de $x$,
le schéma $V$ est régulier en les points de l'image du morphisme
$T→V$. (Un schéma local est régulier si et seulement si son hensélisé
l'est.) En particulier, $V$ est régulier en les points de la
fibre $V_x$. Le lieu régulier étant ouvert, on peut supposer
d'après \ref{2-voisinages fibre restent couvrants} — quitte à
rétrécir le voisinage $U$ de $x$ — que $V$ est régulier.
Enfin, il résulte de \ega{iv}{19.8.1 (ii)} que la propriété d'être
un diviseur à croisements normaux stricts se descend si elle est satisfaite à la limite. 
On conclut par compacité de $X$ pour la topologie étale
complètement décomposée, moins fine que les topologies
$\alt$ et $\altl$. Le cas respé de la proposition \ref{2-réduction cas hensélien}
est un corollaire immédiat de la démonstration du cas non respé.
\end{démo}

\ifx\danslelivre\undefined
\bibliography{../fichiers-configuration-latex/bibliographie}
\bibliographystyle{../fichiers-configuration-latex/smfalpha}
\end{document}
\fi
